% ================================================================
%  18.9 De Profundis (原书 p.926–929)
%  源文: _work/chunks/chunk16.md
% ================================================================
\section{自深处}

\subsection{(1) 感恩}

\notehead{187}{}

(4月7日)这场回顾与盘点——盘点全部《收获与播种》的思考教给了我什么——我相信已经走完
了一遍。唯有一处我从这份回顾里剔除出去,即《收获与播种》的第五部分\fnote{1046}{(*)}{(6月22日)还有第四部分(我眼下正在写的)也一样!见打字稿第1240页脚注(**)。},它至今没有写完。它最初本是“阴阳之钥”里的一段“枝
蔓”,后来延宕了整整一个月,落实为关于 C.G.~Jung自传的一百来页“读书笔记”。由于
这条枝蔓的收尾仍\dmark{}\origpage{1248}遥遥无期,我便把它搁置了。我最盼望的,是把\emph{葬礼}善始善终:写完、打
字、印出、四处寄发——然后,再没有人提它!

何况我预感到,这第五部分将就这同一个\emph{葬礼}给我带来意想不到的光亮——可不是嘛!——
通过我打算在其中对Jung与Sigmund~Freud之间关系的考察:有多年之久,弗洛伊德曾是尚
在寻找道路的年轻Jung的师父。初读自传里论及这层关系的那一章,我什么也没看出来——随后,
若干蹊跷之处引起了我的注意,我折回去看了几处,又把这一章重读了一遍。显然,这层关系绷
满了暧昧——弗洛伊德本人似乎强烈地感到了它,而Jung却津津有味地全然无视它(活脱脱像个
初来的讨论班新生……),把弗洛伊德的不安一概记在他那“神经症”的账上(他还饶有兴致地
把这条“神经症”描绘得有声有色,也许甚至绘声绘色得过了头,反而不大像真的……)。话虽如
此,种种联想随之而来,系于我的朋友、(同时也是)非学生Deligne与我的关系之上;这些联想我
打算追下去,兴许还要稍加开掘。我隐隐觉得,就所牵动的心理机制而言,\emph{葬礼}上发生的事,
绝非一桩绝无仅有、不合典型的巧合——恰恰相反!而我预感到,Jung与弗洛伊德的那层关系,
在这方面没准能提供更多的光亮。

不过对我而言,至少眼下,这第五部分(它也许会叫“Jung——或一场冒险的陷入泥沼”\fnote{1047}{(*)}{本想写“陷入泥沼”(enlisement),却看见自己写出了“葬礼”(enterrement)。这个口误提示的新名字——“Jung——或一场冒险的葬礼”——未必不贴切,甚至没准比原先拟的那个更切中要害。}),已经不再是
那场\emph{葬礼},哪怕它从中生出——我甚至要说:它已不再是《收获与播种》!它是“\textbf{此后}”
(l'Après)——同种种回声一样,那些回声,随着“愚妄与更新”、“葬礼(I)——或中国皇帝的新衣”
和“葬礼(III)——或四种运算”这三部分的寄出,势必要纷纷回到我这里,其中肯定少不了一些
尖酸刻薄的话\fnote{1048}{(**)}{不要把《收获与播种》的第四部分(副题为“四种运算”),与同在该部分之内、集结在这一名目之下
的那组笔记(n\textsuperscript{o}s~167$'$--176$_{7}$)混为一谈。}。到这里,等这部分打完清稿,就已经是一千来页,甚至更多——已经够多的
了!一天的难处一天当就够了……

\pmark\origpage{1249}这种赶紧收场、“把它寄出去”的急切,无疑首先是一匹嗅到火药味的战马的急切,急不可
耐地要冲进战阵\fnote{1049}{(*)}{类似的心绪,在“愚妄与更新”最后一节“一段过去的重量”(n\textsuperscript{o}~50)里已经谈过,只是视角稍有不
同:在那里,冲进战阵的不是战马,而是一头公牛,追赶着一块在它“鼻子跟前晃来晃去”的红布……}。但也许还有更深一层的,是一种愿望:巴望某一段过去从此脱开我而去。这一千来页,以惊
人的方式把这段过去的全部\textbf{重量}凝成了实体——把这项工作做到终了,直到最后一桩庶务
(最后一桩,想必就是把《收获与播种》寄给我寄发清单上预先开列的一百三十位收件人……\fnote{1050}{(**)}{到那时,那著名的“重量”将变得更加“惊人”:一上来就是二十万页(200×1000),而不是一千页!}),在我
看来,几乎出于本能,也正是我把它\textbf{抛掉}的时刻。幻觉?唯有将来知道……

就这样,我来到了那些“终止和弦”面前——在那个有名的“句号”之前;那个“句号”,一年多
来我一直以为它就在眼前,却一天天、一星期一星期、一个月一个月地被推远,总有料想不到的
东西涌来,要占一席之地。

还剩什么要说,在这些终止和弦里?有一份感激,它要借“感谢”说出。这场思考是孤独的果实,
然而我得到过许多方式的帮助。

最显而易见的帮助,来自Zoghman~Mebkhout,同样是以许多方式:凭他的耐心,他领我“入了门”,
走进“好上帝-Mebkhout定理”周身的哲学;凭他对我的信任,不管不顾地把他在我昔日学生们那里
遭遇到的困难与挫折向我和盘托出;凭他帮我厘清一片枝蔓丛生的数学文献——我早已同它断了
往来;最后,凭他自打得知这项工作那一刻起,就对它投以无私而友善的关注——他在其中看见
我倾注了身心,而首先(我相信),他认出并接纳了其中的\textbf{见证}。

我同样感激Pierre~Deligne,他曾专程前来看我,并(去年10月)读阅了当时已写出的\emph{葬礼}
部分,向我提出意见\dmark{}\origpage{1250}\fnote{1051}{(*)}{关于这次造访以及Deligne带给我的种种澄清,见“葬礼(III)”中构成“最后的义务(或造访)”一节的那两条注(n\textsuperscript{o}s~163、164)。}。这次造访对我多有助益。

最后,帮助过我的,还有朗格多克科学技术大学那些为我打字的秘书们的好意与友善氛围:Boulet
小姐(Mlle~Boulet)、Boucher夫人、Brun小姐、Cellier夫人、
Lacan小姐、Mori夫人。她们中有两位,不惜动用个人时间,为按时交
件承担了一部分打字工作,还不肯为此收取报酬——这份情义令我深深感动。此外,整部《收获
与播种》笔记的后一半,是Lacan小姐一个人打出来的,其用心与效率堪称典范。对她们每一位,
我都乐于在此表达我的感激。

我还想着所有那些人,所有那些女人和男人——在我工作的许多时刻,他们曾以常常是不受欢迎
的方式\fnote{1052}{(**)}{这几个月的笔记里,这里那里曾提到这类“打扰”。关于这一点,特别见注“信使(2)”(n\textsuperscript{o}~181)。}打断我的工作与安宁。当然,这些“打扰”本身——它们在某些时刻考验过我,有一些至今还在
我心里留着悲哀的残余——在我的工作里同样有它们的角色要扮演,同样带来一个讯息,听不听、
化不化,全在我自己。当悲哀或怨怼化为感激,我便会知道,这个讯息已被接纳……

\subsection{(2) 朋友}

\notehead{188}{}

\emph{葬礼}的这些终极和弦,近一年来早已名字现成:\emph{自深处}!(De~Profundis)在《引论》
(I~7,“葬仪的安排”)里,我甚至走得更远,竟(也许失之轻率……)预告说,死者那“完全的满
足”将构成“那场难忘葬礼的终曲与最后的和弦”。写下这句预言时(仿佛事已办讫),我是可以
原谅的——写那几行的时候(去年5月),这看上去确实是一个为期极近的预言,因为我自以为正
要抵达“自深处”的那些终止和弦。

诚然,比去年更切身地(当时,思考的“第二口气”行将告一段落),我意识到\dmark{}\origpage{1251}自己离真正“走
完”这葬礼还有多远——若把纯然物质性的事实除外的话(那些事实,我自觉得已经“尽在掌握”,
颇有几分自负\fnote{1053}{(*)}{(5月10日)然而,这几行写下之后,又有一个多月过去,耗在把新冒出来的事实勉力“塞进去”,临
时添了足足二十来条子注!})。倘若确实像我不时感到的那样,理解\emph{葬礼},也就意味着“理解冲
突”,那么很可能,我余生的岁月将不足以完成这番“走完”——至少走不到深处。

如此说来,我可以说,我此刻写下这条终极之注时的心境,已全然不同于写《\emph{葬礼}引论》时
的心境。这是否意味着,我是在“完全的满足”之情不在场的情况下,为这场思考收尾的?

我不以为然。愿景一旦深化,某项曾催生这愿景、为这深化作过铺垫的工作——它本可以被看
作“已臻完成”——便会在一个“彼岸”显现之时,显出自身的\textbf{未完成}。然而,工作的\textbf{意
义},以及它让我们尝到的满足或不满足,并不在它的完成之中,也不取决于这项工作注定完
成与否。工作的意义就在工作本身,就在\textbf{当下此刻}——在我们做事时的心境里,在我们倾
注其中的爱里(或爱的缺席里……)——而不在一个我们无从企及的假想的将来。

去年3月,甚至在我发现\emph{葬礼}之前,我在《引论》(I~1,“梦与实现”,打字稿第~iv~页)里写道:

\begin{quote}
“……我离开这项工作,带着一个知道已把工作做到终点的人的完全满足。没有任何一件事,
无论它多么‘微小’,被我避而不谈;没有任何一件事,是我念念于心要说、却终究没有说的,
以至于在此刻,会在我心里留下一星不满、一星遗憾的残余,无论它们多么‘微小’。”
\end{quote}

如今我知道,这份我原以为“做到终点”的工作,今天仍未到终点,也许永远到不了。但我也知
道,这毕竟是一件次要的事。那份“完全的满足”——在我写下这些力图贴它最近的话语的那一刻,
我曾强烈地感到过它——它一路伴随我写完了\dmark{}\origpage{1252}《收获与播种》。它是我的一位老朋友,早在
我数学生涯的全部岁月里,它就陪伴着我,低声告诉我:路走对了。后来,在冥想的功课里,我
又找到了它——正是同一位。

当我听不见它,工作便失去意义。正因如此,它的声音于我弥足珍贵,也正因如此,我在工作时
小心在意,从不让自己远离它。靠了这个,工作在我一生中始终是喜悦的源泉——那个全然以身
相许者的“完全的满足”。

在即将结束的这项工作里,情形并无二致——这项工作,即是“\textbf{收获}”,同时又是
“\textbf{播种}”。
